\chapter{Lecture 2: New Revolutions in Particle Physics: Standard Model}

This chapter follows Leonard Susskind's second lecture in the supplementary Standard Model sequence, with curation by LazyingArt LLC. The lecture is explicitly modest in scope: one cannot do the full mathematics of quantum chromodynamics in a single meeting, so the aim is to isolate the basic picture. The route is deliberate. We first place QCD beside QED, then step backward into the mathematics of spin, and only after that use the same pattern of state-counting to reach isospin, color, gluons, and the first qualitative picture of QCD.

\section{QCD by Analogy, and Why We Must Revisit Spin}

Quantum electrodynamics is the theory of electrons and photons, and in the atomic setting it is also the theory of the Coulomb force that binds matter together. Susskind begins by asking us to think of quantum chromodynamics in the same spirit. QCD is the theory of quarks and gluons, and also of the composite objects they make. Those composite objects are hadrons. Among them are baryons, such as the proton and neutron, which carry net baryon number and in this lecture are represented schematically as three-quark states, and mesons, which are quark--antiquark states. In both cases the gluons are the binding medium, the analogue of the field that binds electrically charged matter.

But before turning to quarks and gluons directly, the lecture insists on a return to spin. This is not a digression. The mathematics of angular momentum will immediately reappear in another guise as isospin, and a closely related pattern of state-counting will later organize color. The lecture does not want a formal excursion into group theory; instead it uses the familiar language of angular momentum as the working substitute.

\section{Angular Momentum Multiplets and Two Spin-Half Particles}

The starting point is the algebra of angular momentum:
\begin{equation}
[L_x,L_y]=i\hbar L_z,
\end{equation}
together with the two cyclic companions
\begin{equation}
[L_y,L_z]=i\hbar L_x,
\qquad
[L_z,L_x]=i\hbar L_y.
\end{equation}
The physics is rotationally symmetric, but in solving the problem we choose a preferred axis, conventionally the \(z\)-axis. This is not symmetry breaking in nature; it is only a choice in our description. Since the components do not commute, we cannot simultaneously specify all three of them. So we focus on \(L_z\), and the standard angular-momentum machinery gives a multiplet of states labeled by a maximum value \(\ell\), with
\begin{equation}
m_z=\ell,\ell-1,\ldots,-\ell,
\qquad
\#\text{ of states}=2\ell+1.
\end{equation}
The lecture says \(L\) aloud for the multiplet label, while the board writes \(\ell\); here we use \(\ell\) for the multiplet and reserve \(L_x,L_y,L_z\) for the operators.

For a single spin-half particle there are two states. For a spin-one object there are three. The useful point is not just the list of possibilities but the counting rule itself: once the maximum value is known, the size of the multiplet is fixed.

Now let us combine two spin-half particles and ignore orbital angular momentum. The product basis is
\begin{equation}
|\uparrow\uparrow\rangle,\qquad
|\uparrow\downarrow\rangle,\qquad
|\downarrow\uparrow\rangle,\qquad
|\downarrow\downarrow\rangle.
\end{equation}
The maximum \(z\)-component is \(+1\), attained by \( |\uparrow\uparrow\rangle \), and the minimum is \(-1\), attained by \( |\downarrow\downarrow\rangle \). But there are two distinct states with \(m_z=0\), namely \( |\uparrow\downarrow\rangle \) and \( |\downarrow\uparrow\rangle \). That is where the first real tension in the lecture appears.

\begin{figure}[t]
\centering
\includegraphics[width=0.78\textwidth]{lecture_02_figure_02.png}
\caption{Symmetric mixed-spin state highlighted on the board. The circled sum is the state singled out before the normalization is written explicitly.}
\end{figure}

The board at this stage isolates the symmetric mixed-spin combination. In cleaned form we write it as
\begin{equation}
\frac{|\uparrow\downarrow\rangle+|\downarrow\uparrow\rangle}{\sqrt2}.
\end{equation}
The normalization factor is not yet visible in the earlier board frame, but it is required once we pass from a circled sum to a normalized quantum state.

\subsection{Question \& Answer}

\paragraph{Question.} If two spin-half particles give four basis states, why do they not form a single spin-one multiplet?

\paragraph{Answer.} Because a spin-one multiplet contains only three states:
\[
m_z=+1,\ 0,\ -1.
\]
The two aligned states already account for the \(m_z=\pm1\) members. A spin-one multiplet can therefore accommodate only one linear combination of the two \(m_z=0\) basis states. The second independent \(m_z=0\) combination must lie outside the spin-one multiplet, and the only remaining possibility is spin zero.

\section{Triplet, Singlet, and the Meaning of the Symmetric State}

The lecture now resolves the counting puzzle by diagonalizing the two-dimensional \(m_z=0\) sector into symmetric and antisymmetric combinations. The result is the familiar but here highly motivated decomposition
\begin{align}
\frac12\otimes\frac12 &= 1\oplus 0, \\
\ell=1:&\qquad
|\uparrow\uparrow\rangle,\qquad
\frac{|\uparrow\downarrow\rangle+|\downarrow\uparrow\rangle}{\sqrt2},\qquad
|\downarrow\downarrow\rangle, \\
\ell=0:&\qquad
\frac{|\uparrow\downarrow\rangle-|\downarrow\uparrow\rangle}{\sqrt2}.
\end{align}

\begin{figure}[t]
\centering
\includegraphics[width=0.80\textwidth]{lecture_02_figure_03.png}
\caption{Triplet and singlet decomposition of two spin-\(\tfrac12\) states. The board organization itself is important: the three \(\ell=1\) states are grouped together, while the antisymmetric \(\ell=0\) state is written separately below.}
\end{figure}

The lecture does not leave the matter at the level of algebraic bookkeeping. It insists on the physical meaning of the decomposition. The symmetric mixed-spin state has \(m_z=0\), but it does \emph{not} have total angular momentum zero. It belongs to the \(\ell=1\) triplet. This is the state that allows the spin-one multiplet to rotate into other directions. Zero along the chosen \(z\)-axis is not the same thing as zero total angular momentum.

The antisymmetric combination is different:
\begin{equation}
\frac{|\uparrow\downarrow\rangle-|\downarrow\uparrow\rangle}{\sqrt2}
\quad \longleftrightarrow \quad
\ell=0.
\end{equation}
As far as angular momentum is concerned, it behaves like nothing at all.

This is sharpened by the lecture's deliberately provocative question: if we worried only about angular-momentum conservation, which two-spin combination could simply disappear? The answer is not either of the aligned states, and not the symmetric \(m_z=0\) state either, because those all belong to \(\ell=1\). Only the antisymmetric singlet could do so. In that sense the singlet is the true zero-angular-momentum state.

There is a further subtle point that the lecture emphasizes in response to a question from the audience. If we measure the symmetric mixed-spin state along the \(z\)-axis, we get zero. But that same state is not inert under rotations. Measure along another axis, and it participates in the same spin-one structure as the \(m_z=\pm1\) states. This is exactly why the decomposition into triplet and singlet matters.

\section{Three Spins as a Template for Isospin}

Before turning to hadrons, Susskind briefly repeats the same counting logic for three spin-half objects. There are
\begin{equation}
2\times2\times2 = 8
\end{equation}
product states altogether. The maximum \(z\)-component is \(+\tfrac32\), so a spin-\(\tfrac32\) multiplet must appear, and it contains
\begin{equation}
2\left(\frac32\right)+1 = 4
\end{equation}
states. Four states therefore remain. They cannot form another spin-\(\tfrac32\) multiplet, because the \(\pm\tfrac32\) extrema have already been used up. The only remaining possibility is spin-\(\tfrac12\), and since four states remain, that representation occurs twice:
\begin{equation}
\frac12\otimes\frac12\otimes\frac12
=
\frac32\oplus\frac12\oplus\frac12.
\end{equation}

The lecture then asks why we are doing this. The answer is that the same combinatorics will organize a new internal symmetry, isospin. Isospin is not an actual rotation of physical space. It is an internal two-state structure whose mathematics is isomorphic to the spin-half system we have just reviewed.

Why is this a good approximation? Because the natural scale of hadronic physics is a few hundred MeV, while the up and down quarks are both very light on that scale and, more importantly, close in mass to one another. If we also ignore electromagnetic effects at first pass, then the strong interaction sees \(u\) and \(d\) as nearly interchangeable. In that approximation the pair
\[
u,\qquad d
\]
behaves like a spin-half doublet. Historically isospin came first and quarks later; in this lecture, however, it is introduced as the mathematical language that organizes the light-flavor sector.

\section{From Spin to Isospin: Proton and Neutron}

A single quark is not treated in the lecture as an isolated laboratory object. The first experimentally accessible hadrons we discuss are the proton and neutron, each built from three quarks. By the same counting logic that we used for three ordinary spins, three light quarks can organize themselves into an isospin-\(\tfrac12\) sector and an isospin-\(\tfrac32\) sector. We start with the gentler case:
\begin{equation}
p\sim uud,
\qquad
n\sim udd.
\end{equation}
These are the two states of the isospin-\(\tfrac12\) doublet.

Their electric charges follow from
\begin{equation}
Q(u)=+\frac23,
\qquad
Q(d)=-\frac13.
\end{equation}
Thus
\[
Q(p)=\frac23+\frac23-\frac13=+1,
\qquad
Q(n)=\frac23-\frac13-\frac13=0.
\]
The near equality of their masses is one of the lecture's concrete signs that the interchange \(u\leftrightarrow d\) is a good approximate symmetry of the strong interaction.

To make the fermionic issue precise, the lecture temporarily labels the three quark slots by \(1,2,3\). The board then writes an antisymmetrized factor involving the two up quarks:
\begin{equation}
d_1\bigl(u_2u_3-u_3u_2\bigr).
\end{equation}
This is the exact blackboard expression we retain.

\begin{figure}[t]
\centering
\includegraphics[width=0.58\textwidth]{lecture_02_figure_04.png}
\caption{Antisymmetrized two-\(u\)-quark factor in the proton state. The board records only this factor, not the full proton wavefunction.}
\end{figure}

A caution is essential here. The lecture is not constructing the full proton wavefunction on the board; it is isolating the antisymmetric factor that plays the role of the singlet-like piece in the discussion. If one wished to normalize that two-\(u\)-quark combination in standard notation, one would insert a factor of \(1/\sqrt2\), but that factor is not written on the board and is not needed for the logic of the lecture.

The proton and neutron are therefore introduced as the hadronic analogue of the spin-half doublet: two states, exchanged by \(u\leftrightarrow d\), and nearly degenerate once we ignore small mass differences and electromagnetic effects.

\section{The Delta Multiplet and the Pauli Puzzle}

The lecture then turns to the more dramatic case: the isospin-\(\tfrac32\) multiplet. By analogy with the spin-\(\tfrac32\) quartet, we expect four flavor states,
\begin{equation}
\Delta^{++}\sim uuu,\qquad
\Delta^{+}\sim uud,\qquad
\Delta^{0}\sim udd,\qquad
\Delta^{-}\sim ddd.
\end{equation}
These are not the proton and neutron again. The crucial difference is the ordinary spin: the \(\Delta\) states are spin-\(\tfrac32\), not spin-\(\tfrac12\). In the lecture this means that the three ordinary spins are aligned, while the isospins are aligned as well.

Their charges are
\begin{align}
Q(\Delta^{++})&=+2, &
Q(\Delta^{+})&=+1, &
Q(\Delta^{0})&=0, &
Q(\Delta^{-})&=-1.
\end{align}
Within the approximation of isospin symmetry, the four members of the quartet are nearly degenerate. The lecture quotes the proton and neutron masses as about \(940\) MeV, while the \(\Delta\) mass is about \(1200\) MeV.

A useful worked example is the decay of the doubly charged member. The lecture sketches it in quark language by inserting a \(d\bar d\) pair:
\begin{align}
\Delta^{++} &\sim uuu, \\
uuu &\longrightarrow uud + u\bar d \\
&= p + \pi^+.
\end{align}
The charge bookkeeping is immediate:
\[
Q(p)=+1,
\qquad
Q(\pi^+)=Q(u)+Q(\bar d)=\frac23+\frac13=+1,
\]
so the final state again has total charge \(+2\). The decay is energetically allowed because
\[
1200\ \text{MeV} \; > \; 940\ \text{MeV} + 140\ \text{MeV},
\]
leaving room for kinetic energy. This is why the \(\Delta\) resonances are real but short-lived hadronic states.

\subsection{Question \& Answer}

\paragraph{Question.} How can a \(\Delta^{++}\) made of three up quarks avoid violating Fermi statistics?

\paragraph{Answer.} If spin and flavor were the whole story, it could not. In \(\Delta^{++}\) the three quarks are all \(u\), and because the baryon has spin \(\tfrac32\), their ordinary spins are aligned as well. The state would then appear to contain three identical fermions in the same state, which is forbidden by the Pauli principle. The lecture treats this as the historical clue that forces a new hidden label into the theory. Quarks must carry one more quantum number so that the three quarks in the \(\Delta^{++}\) are not actually identical in all respects. That extra label is color.

\section{Color, Gluons, and the Minimal Picture of QCD}

Color is introduced as a three-valued label, conventionally
\[
r,\qquad g,\qquad b.
\]
Every flavor of quark comes in these three versions, and antiquarks carry the corresponding anti-colors
\[
\bar r,\qquad \bar g,\qquad \bar b.
\]
The lecture insists that these are only labels. They are not colors in the everyday sense, and they do not differ in mass or any other obvious property. Their role is to distinguish states that would otherwise violate Fermi statistics.

In the lecture's matrix-style bookkeeping, a quark can be organized as a three-component column,
\begin{equation}
q=
\begin{pmatrix}
r\\
g\\
b
\end{pmatrix},
\qquad
\bar q=
\begin{pmatrix}
\bar r & \bar g & \bar b
\end{pmatrix},
\end{equation}
so that a gluon is temporarily thought of as a color--anticolor object filling out a \(3\times 3\) matrix:
\begin{equation}
q\,\bar q \sim
\begin{pmatrix}
r\bar r & r\bar g & r\bar b\\
g\bar r & g\bar g & g\bar b\\
b\bar r & b\bar g & b\bar b
\end{pmatrix}.
\end{equation}
This gives nine possibilities at first sight. The lecture then adds the important caveat that one diagonal superposition is not an independent physical gluon. A minimal cleaned completion of that statement is
\begin{equation}
3\otimes \bar 3 = 8\oplus 1.
\end{equation}
The full reason is deferred; the lecture only wants the audience to register the ``nine minus one'' subtlety.

The basic QCD vertex is a color-changing emission process. In cleaned notation we may write
\begin{equation}
q^a \longrightarrow q^b + g^{a\bar b}.
\end{equation}
For example, a red quark can become a green quark by emitting a red--antigreen gluon:
\begin{equation}
r \longrightarrow g + (r\bar g).
\end{equation}
The lecture's rule for working with these diagrams is operational rather than formal: follow the color lines. When a line turns around in time, reinterpret it as an antiparticle.

\begin{figure}[t]
\centering
\begin{tikzpicture}[scale=1.0]
\draw[->] (0,0)--(2,0) node[midway,above] {$r$};
\draw[->] (2,0)--(4,1) node[midway,above] {$g$};
\draw[dashed] (2,0)--(4,-1) node[midway,below] {$r\bar g$};
\fill (2,0) circle (1.5pt);
\end{tikzpicture}

\bigskip

\begin{tikzpicture}[scale=1.0]
\draw[dashed] (0,0)--(2,0) node[midway,above] {$r\bar b$};
\draw[dashed] (2,0)--(4,1) node[midway,above] {$r\bar g$};
\draw[dashed] (2,0)--(4,-1) node[midway,below] {$g\bar b$};
\fill (2,0) circle (1.5pt);
\end{tikzpicture}
\caption{Schematic QCD vertices in the spirit of the lecture. Above: a quark changes color and emits a gluon. Below: a gluon turns into two gluons. The labels are bookkeeping labels, and the governing rule is to follow the color lines.}
\end{figure}

The upper diagram is the analogue of photon emission in QED. The lower one is the crucial novelty. Photons do not carry electric charge, so at the primitive level they do not interact with one another. Gluons are different. Because the gluon itself carries color bookkeeping labels, gluons can interact with gluons. The lecture treats this as the decisive structural difference between QCD and QED.

That immediately changes the dynamics. Gluons can be exchanged between quarks and create forces between quarks, just as photons do between electrically charged particles. But gluons can also be exchanged between gluons. This means that QCD contains forces between its own force carriers. In the lecture this leads directly to three qualitative consequences:

First, the gluon dynamics is non-linear. A gluon wave does not simply pass through another gluon wave the way two beams of light do in vacuum. Second, one can have bound objects in QCD made only of gluonic degrees of freedom. Third, the bookkeeping of color flow is more elaborate than in QED, which is why the lecture keeps returning to the rule ``follow the lines'' rather than trying to encode every example in abstract notation.

The lecture closes with a heuristic rather than a derivation. Imagine a small mass \(m\) attached to a soft, extended piece of matter. Push very slowly, and the whole assembly moves together; the effective inertia is the mass of the full composite. Shake very rapidly, and the surrounding material does not have time to respond; now one probes only the core. The observed mass depends on the frequency of the probe. This is the lecture's route to the idea that the parameters of a quark depend on scale:
\begin{equation}
m_q=m_q(\mu),
\qquad
Q_q=Q_q(\mu).
\end{equation}
These are the running constants of the theory in the lecture's qualitative sense. No renormalization-group derivation is attempted here. The point is only that ``the mass of a quark'' is not a single rigid number independent of how the quark is probed.

The deepest question is postponed. Why do we never isolate a single quark in the laboratory? The lecture names the issue---confinement---but defers its real explanation to the next meeting.

\section{Summary}

The chapter has followed the lecture's intended arc. QCD was introduced first by analogy with QED, but the real work began with angular-momentum multiplets. The two-spin counting puzzle led to the triplet and singlet, the three-spin counting led to the \(\tfrac32\) and \(\tfrac12\) sectors, and that same mathematics then reappeared as isospin for the light quarks. From there the proton and neutron emerged as an isospin doublet, the \(\Delta\) quartet sharpened the Pauli puzzle, and color entered as the hidden label that resolves it.

Only after that historical and conceptual pressure had built up did the lecture turn to QCD proper: three colors, gluons organized as color--anticolor bookkeeping states, quark--gluon emission, gluon self-interaction, and the resulting non-linear dynamics. The lecture ends where it intends to end: not with a complete theory of hadrons, but with an intuitive picture and a clear promise that confinement and hadron structure come next.